% ==============================================================================
% Chapter 10: RNA Sequencing — Measuring Gene Expression
% ==============================================================================
\chapter{RNA Sequencing: Measuring Gene Expression}
\label{ch:rnaseq}

\begin{keyconcept}
RNA sequencing (\acrshort{RNA-seq}) measures the transcriptome---the complete set of RNA
transcripts in a cell, tissue, or organism at a given moment. Because gene
expression is the primary mechanism through which genotype becomes phenotype,
\acrshort{RNA-seq} has become one of the most widely used sequencing applications in
biology and medicine.
\end{keyconcept}

% ---------------------------------------------------------------------------
\section{Why Measure RNA?}
\label{sec:rnaseq:why}

Every cell in a multicellular organism carries essentially the same genome, yet
a neuron looks and behaves nothing like a hepatocyte. The difference lies
predominantly in \emph{which genes are expressed} and at \emph{what levels}.
Measuring RNA therefore provides a direct readout of cellular state.

\subsection{Gene Expression as a Window into Biology}

Gene expression governs virtually every biological process:

\begin{itemize}[nosep]
  \item \textbf{Development.} Transcription factors activate stage-specific
    gene programs that drive differentiation from a single zygote into hundreds
    of specialized cell types.
  \item \textbf{Disease.} Cancers are defined by aberrant expression profiles;
    autoimmune diseases by dysregulated immune genes; neurodegeneration by
    altered synaptic gene networks.
  \item \textbf{Environmental response.} Cells respond to stress, infection,
    drugs, and nutrients by rapidly reprogramming their transcriptomes.
  \item \textbf{Circadian rhythm.} Thousands of genes oscillate with 24-hour
    periodicity in virtually every tissue.
\end{itemize}

Before \acrshort{RNA-seq}, gene expression was measured by microarrays,
Northern blots, or quantitative PCR (qPCR). Microarrays could interrogate
thousands of genes simultaneously but relied on predefined probes, suffered
from cross-hybridization, and had limited dynamic range. \acrshort{RNA-seq}
overcame all of these limitations: it is unbiased, quantitative across
five orders of magnitude, and can detect novel transcripts, splice variants,
and allele-specific expression without any prior knowledge of the genome.

\subsection{The Transcriptome}

The transcriptome is far more complex than a simple catalogue of protein-coding
messenger RNAs:

\begin{description}[style=nextline]
  \item[Messenger RNA (mRNA)] Protein-coding transcripts, accounting for only
    $\sim$2--5\% of total RNA mass but $\sim$20,000 genes in humans.
  \item[Ribosomal RNA (rRNA)] The dominant RNA species by mass ($\sim$80--85\%),
    forming the structural and catalytic core of ribosomes (28S, 18S, 5.8S, 5S
    in eukaryotes).
  \item[Transfer RNA (tRNA)] Adapter molecules ($\sim$10--15\% of RNA) that
    deliver amino acids during translation.
  \item[Long non-coding RNA (lncRNA)] Transcripts $>$200\,nt that do not encode
    proteins; $>$60,000 annotated in humans. Functions include chromatin
    remodeling, transcriptional regulation, and scaffolding (e.g., \gene{XIST},
    \gene{MALAT1}, \gene{HOTAIR}).
  \item[Small non-coding RNAs] Include microRNAs (miRNAs, $\sim$22\,nt),
    piwi-interacting RNAs (piRNAs, 24--32\,nt), small interfering RNAs
    (siRNAs), small nucleolar RNAs (snoRNAs), and tRNA-derived fragments
    (tRFs).
  \item[Circular RNAs (circRNAs)] Covalently closed loops formed by
    back-splicing; highly stable, tissue-specific, and implicated in miRNA
    sponging and translation regulation.
  \item[Transcript isoforms] Alternative splicing, alternative polyadenylation,
    and alternative transcription start sites generate multiple isoforms per
    gene. Human genes produce an average of $\sim$7 distinct isoforms.
\end{description}

\begin{warningbox}
The vast majority of cellular RNA is ribosomal RNA. Without an enrichment or
depletion strategy, an \acrshort{RNA-seq} experiment would waste $>$80\% of
sequencing reads on rRNA.
\end{warningbox}


% ---------------------------------------------------------------------------
\section{Bulk RNA-seq: An Overview of Approaches}
\label{sec:rnaseq:bulk}

Bulk \acrshort{RNA-seq} measures the \emph{average} expression across a
population of cells (typically $10^5$--$10^7$ cells). Two main strategies exist,
distinguished by how they deal with the rRNA problem.

\subsection{mRNA-seq: Poly-A Selection}

The most common approach for studying protein-coding gene expression.
Polyadenylated transcripts (mRNAs plus some lncRNAs) are captured using
oligo-dT beads that hybridize to the poly-A tail.

\begin{protocolbox}[title={Poly-A Selection Protocol Overview}]
\begin{enumerate}[nosep]
  \item Extract total RNA (e.g., TRIzol or column-based methods).
  \item Incubate with oligo-dT magnetic beads; poly-A$^+$ transcripts bind.
  \item Wash away rRNA, tRNA, and other non-polyadenylated RNAs.
  \item Elute mRNA and proceed to fragmentation and cDNA synthesis.
\end{enumerate}
Requires high-quality RNA with intact poly-A tails (RIN $\geq$ 7 recommended).
\end{protocolbox}

\textbf{Advantages:}
\begin{itemize}[nosep]
  \item Highly enriched for protein-coding transcripts.
  \item Lower sequencing depth required (fewer ``wasted'' reads).
  \item Well-established, widely used---most public datasets use this approach.
\end{itemize}

\textbf{Limitations:}
\begin{itemize}[nosep]
  \item Misses non-polyadenylated transcripts (histone mRNAs, many lncRNAs, most
    small RNAs, bacterial transcripts).
  \item Requires high-quality RNA; degraded samples lose poly-A tails, biasing
    toward 3$'$ ends.
  \item Cannot capture nascent or nuclear RNA species lacking poly-A tails.
\end{itemize}

\subsection{Total RNA-seq: Ribosomal RNA Depletion}

Ribosomal depletion removes rRNA from total RNA, leaving all other RNA species
intact---mRNAs, lncRNAs, circRNAs, pre-mRNAs, and other non-coding species.

Common rRNA depletion kits include:
\begin{itemize}[nosep]
  \item \textbf{Illumina Ribo-Zero Plus:} uses DNA probe hybridization and
    RNase H digestion to deplete rRNA from human, mouse, rat, and bacterial
    samples.
  \item \textbf{RiboCop (Lexogen):} enzymatic rRNA depletion compatible with
    degraded samples.
  \item \textbf{NEBNext rRNA Depletion Kit:} probe-based depletion for human,
    mouse, and rat.
  \item \textbf{QIAseq FastSelect:} targets rRNA during reverse transcription,
    eliminating a separate depletion step.
\end{itemize}

\textbf{Advantages:}
\begin{itemize}[nosep]
  \item Captures both coding and non-coding RNA.
  \item Works with degraded/FFPE samples (no reliance on poly-A tail integrity).
  \item Detects intronic reads, pre-mRNAs, and nascent transcription signals.
  \item Can reveal circular RNAs (which lack poly-A tails).
\end{itemize}

\textbf{Limitations:}
\begin{itemize}[nosep]
  \item Residual rRNA ($\sim$1--10\%) consumes some sequencing capacity.
  \item More complex read composition; requires deeper sequencing for equivalent
    mRNA coverage.
  \item May require organism-specific depletion probes for non-model organisms.
\end{itemize}

\subsection{Choosing Between mRNA-seq and Total RNA-seq}

\begin{table}[htbp]
\centering
\caption{Comparison of mRNA-seq (poly-A selection) and total RNA-seq (rRNA
depletion) approaches.}
\label{tab:rnaseq:mrna_vs_total}
\small
\begin{tabularx}{\textwidth}{lXX}
\toprule
\textbf{Feature} & \textbf{mRNA-seq (poly-A)} & \textbf{Total RNA-seq (rRNA depletion)} \\
\midrule
Target RNA & Poly-A$^+$ mRNA (+ some lncRNA) & All RNA except rRNA \\
RNA quality requirement & High (RIN $\geq$ 7) & Moderate (works with RIN $\geq$ 2) \\
FFPE samples & Poor performance & Suitable \\
Non-coding RNA & Missed (most) & Captured \\
Intronic reads & Minimal & Substantial \\
Circular RNA & Not captured & Captured \\
Typical depth needed & 20--30M reads & 40--60M reads \\
Cost per informative read & Lower & Higher \\
Best for & Standard differential expression & Degraded samples, ncRNA, comprehensive profiling \\
\bottomrule
\end{tabularx}
\end{table}

\begin{tipbox}
If your primary goal is differential expression of protein-coding genes and your
RNA quality is good (RIN $\geq$ 7), mRNA-seq is the most cost-effective choice.
If you are working with clinical/FFPE samples, interested in non-coding RNAs, or
need to capture the broadest possible transcriptome, use total RNA-seq.
\end{tipbox}


% ---------------------------------------------------------------------------
\section{Library Preparation for RNA-seq}
\label{sec:rnaseq:library}

\subsection{RNA Extraction and Quality Assessment}

High-quality RNA is the single most critical determinant of a successful
\acrshort{RNA-seq} experiment. RNA is chemically fragile and ubiquitous RNases
in the environment degrade it rapidly.

\begin{description}[style=nextline]
  \item[Extraction methods]
    \begin{itemize}[nosep]
      \item \textbf{TRIzol / TRI Reagent:} acid guanidinium--phenol--chloroform
        extraction. Gold standard for yield; suitable for most tissues.
      \item \textbf{Column-based:} RNeasy (Qiagen), PureLink (Invitrogen). Faster,
        less toxic, but lower yield for difficult tissues.
      \item \textbf{Automated:} KingFisher (Thermo), QIAsymphony. Essential for
        high-throughput core facilities.
    \end{itemize}
  \item[Quality metrics]
    \begin{itemize}[nosep]
      \item \textbf{RIN (RNA Integrity Number):} Agilent Bioanalyzer / TapeStation
        score from 1 (fully degraded) to 10 (intact). Based on the 28S/18S rRNA
        ratio and the electropherogram trace. RIN $\geq$ 7 is recommended for
        poly-A selection.
      \item \textbf{DV200:} Fraction of RNA fragments $>$ 200 nucleotides.
        Preferred metric for degraded/FFPE RNA because the RIN algorithm
        performs poorly when rRNA peaks are absent. DV200 $\geq$ 30\% is
        often the minimum for library construction with specialized FFPE kits.
      \item \textbf{Concentration:} Qubit fluorometer (RNA HS Assay) provides
        accurate quantitation; NanoDrop UV absorbance gives purity ratios
        (260/280 $\sim$ 2.0 for pure RNA; 260/230 $\geq$ 1.8).
    \end{itemize}
\end{description}

\subsection{Strand-Specific Library Preparation}
\label{sec:rnaseq:stranded}

In unstranded \acrshort{RNA-seq}, the original strand of the RNA transcript
cannot be determined from the sequencing reads. Strand-specific (stranded)
protocols preserve this information.

\subsubsection{Why Strandedness Matters}

\begin{itemize}[nosep]
  \item \textbf{Overlapping genes:} In complex genomes, genes on opposite
    strands can overlap. Unstranded data cannot assign reads to the correct
    gene.
  \item \textbf{Antisense transcription:} Many loci produce antisense
    transcripts that regulate the sense gene. Without strand information, sense
    and antisense reads are indistinguishable.
  \item \textbf{Accurate quantification:} Stranded data improves the accuracy
    of gene-level and transcript-level quantification by 5--15\% for
    overlapping annotations.
  \item \textbf{Novel transcript discovery:} Strand information is essential
    for correctly assembling and annotating novel transcripts.
\end{itemize}

\subsubsection{The dUTP Method}

The most widely used strand-specific protocol:

\begin{enumerate}[nosep]
  \item Fragment mRNA and synthesize first-strand cDNA using random primers.
  \item Synthesize second-strand cDNA incorporating dUTP instead of dTTP.
  \item Ligate sequencing adapters.
  \item Digest the dUTP-containing second strand with USER enzyme (uracil-DNA
    glycosylase + endonuclease VIII).
  \item Only the first-strand cDNA (complementary to the original mRNA) is
    amplified and sequenced.
  \item Read 2 corresponds to the sense strand of the original mRNA.
\end{enumerate}

This is the basis of the \textbf{Illumina Stranded mRNA Prep} and
\textbf{Illumina Stranded Total RNA Prep} kits, which have largely become the
community standard.

% TikZ diagram: Strand-specific library prep
\begin{figure}[htbp]
\centering
\begin{tikzpicture}[
  node distance=0.7cm,
  every node/.style={font=\small},
  rna/.style={draw, thick, fill=red!15, rounded corners, minimum width=6cm, minimum height=0.6cm},
  cdna1/.style={draw, thick, fill=blue!15, rounded corners, minimum width=6cm, minimum height=0.6cm},
  cdna2/.style={draw, thick, fill=orange!15, rounded corners, minimum width=6cm, minimum height=0.6cm, densely dashed},
  arrow/.style={-{Stealth[length=6pt]}, thick},
  label/.style={font=\small\itshape, text=gray!70!black}
]

% Step 1: mRNA
\node[rna] (mrna) {5$'$ ------- mRNA (sense) ------- AAAAAAA 3$'$};
\node[label, right=0.3cm of mrna] {Original transcript};

% Step 2: First-strand cDNA
\node[cdna1, below=1.0cm of mrna] (ss1) {3$'$ ------- 1st strand cDNA ------- 5$'$};
\node[label, right=0.3cm of ss1] {Reverse transcription};
\draw[arrow] (mrna.south) -- node[right, font=\scriptsize] {RT + random primers} (ss1.north);

% Step 3: Second-strand with dUTP
\node[cdna2, below=1.0cm of ss1] (ss2) {5$'$ - - - 2nd strand cDNA (dUTP) - - - 3$'$};
\node[label, right=0.3cm of ss2] {Contains dUTP};
\draw[arrow] (ss1.south) -- node[right, font=\scriptsize] {2nd strand synthesis (dUTP)} (ss2.north);

% Step 4: USER enzyme digestion
\node[below=1.0cm of ss2, font=\small, text=BrickRed] (digest) {USER enzyme digests dUTP strand};
\draw[arrow, BrickRed] (ss2.south) -- (digest.north);

% Step 5: Only first strand remains
\node[cdna1, below=1.0cm of digest] (final) {3$'$ ------- 1st strand cDNA ------- 5$'$};
\node[label, right=0.3cm of final] {Amplified \& sequenced};
\draw[arrow] (digest.south) -- (final.north);

% Annotation
\node[below=0.5cm of final, font=\footnotesize\itshape, text=MidnightBlue] {Read 1 $\rightarrow$ antisense; \quad Read 2 $\rightarrow$ sense (original mRNA direction)};

\end{tikzpicture}
\caption{Schematic of the dUTP strand-specific library preparation method. The
second-strand cDNA is synthesized with dUTP (dashed orange) and subsequently
digested by USER enzyme, ensuring that only the first-strand cDNA (blue) is
amplified. This preserves strand-of-origin information.}
\label{fig:rnaseq:stranded}
\end{figure}

\subsection{Low-Input RNA-seq}

Standard RNA-seq protocols require 100\,ng--1\,$\mu$g of total RNA. Many
biological questions demand profiling from small cell populations, sorted cells,
or limited clinical material.

\begin{description}[style=nextline]
  \item[Smart-seq2] Full-length transcript coverage from single cells or
    picogram-level input. Uses template-switching oligos and limited PCR
    amplification. Plate-based (typically 96- or 384-well). Excellent
    sensitivity but low throughput and not strand-specific in the original
    protocol.
  \item[SMART-Seq Stranded (Takara)] Combines Smart-seq template-switching
    with strand-specific library prep. Input as low as 10\,pg (single cells)
    to 10\,ng. Adds strand specificity lacking in original Smart-seq2.
  \item[NEBNext Single Cell / Low Input RNA Library Prep] Template-switching
    approach compatible with 2\,pg--200\,ng input.
\end{description}

\subsection{Degraded RNA and FFPE Samples}

Formalin-fixed paraffin-embedded (FFPE) tissues are ubiquitous in pathology
archives. The fixation and embedding process crosslinks and fragments RNA
severely (RIN typically 1--3). Specialized protocols are needed:

\begin{itemize}[nosep]
  \item \textbf{Illumina Stranded Total RNA with Ribo-Zero Plus:} works with
    FFPE RNA down to DV200 of $\sim$30\%.
  \item \textbf{KAPA RNA HyperPrep with RiboErase:} optimized for low-quality
    input; produces libraries from as little as 25\,ng FFPE RNA.
  \item \textbf{Takara SMARTer Stranded Total RNA-Seq Kit v3 -- Pico Input:}
    designed for 250\,pg--10\,ng of degraded RNA.
  \item \textbf{Probe-based approaches (Twist Bioscience RNA Panel, IDT xGen
    RNA):} use hybridization capture for targeted RNA panels from FFPE
    samples.
\end{itemize}

\begin{warningbox}
FFPE RNA-seq data has distinct characteristics: shorter insert sizes, higher
duplicate rates, and 3$'$ bias. Specialized QC metrics (DV200, percentage of
reads in gene bodies) should be used. The standard RIN score is unreliable for
FFPE samples.
\end{warningbox}


% ---------------------------------------------------------------------------
\section{Small RNA Sequencing}
\label{sec:rnaseq:smallrna}

Small RNA-seq targets RNA molecules $<$ 200\,nt, with a particular focus on
microRNAs (miRNAs, $\sim$22\,nt), which are key post-transcriptional
regulators of gene expression.

\subsection{Small RNA Species}

\begin{description}[style=nextline]
  \item[microRNAs (miRNAs)] 19--25\,nt single-stranded RNAs that bind the
    3$'$ UTR of target mRNAs to induce translational repression or mRNA
    degradation. $\sim$2,600 mature miRNAs annotated in humans (miRBase v22.1).
  \item[piRNAs] 24--32\,nt RNAs that silence transposable elements in the
    germline; associate with PIWI-clade Argonaute proteins.
  \item[siRNAs] 20--24\,nt double-stranded RNA-derived species involved in RNA
    interference.
  \item[tRNA-derived fragments (tRFs)] 14--35\,nt fragments cleaved from mature
    or precursor tRNAs. Emerging roles in stress response, gene regulation,
    and cancer.
  \item[snoRNAs] 60--300\,nt; guide chemical modifications of rRNA and snRNA.
\end{description}

\subsection{Small RNA Library Preparation}

\begin{protocolbox}[title={Small RNA-seq Library Preparation}]
\begin{enumerate}[nosep]
  \item \textbf{Size selection:} Total RNA is size-selected for 15--50\,nt
    fragments using PAGE gel or bead-based selection.
  \item \textbf{3$'$ adapter ligation:} An adapter is ligated to the 3$'$-OH of
    small RNAs (which have a free 3$'$-OH after Dicer processing).
  \item \textbf{5$'$ adapter ligation:} An adapter is ligated to the 5$'$
    phosphate.
  \item \textbf{Reverse transcription:} cDNA is synthesized from the
    adapter-ligated small RNA.
  \item \textbf{PCR amplification:} Limited-cycle PCR with indexed primers.
  \item \textbf{Size selection:} Final library is size-selected ($\sim$135--170\,bp
    including adapters) to remove adapter dimers.
\end{enumerate}
\end{protocolbox}

\begin{warningbox}[title={Adapter Dimers in Small RNA-seq}]
Adapter dimers (3$'$ and 5$'$ adapters ligated directly to each other without
a small RNA insert) are the most common artifact in small RNA-seq. They are
$\sim$120--130\,bp and can dominate the library if input RNA is limiting.
Rigorous gel-based or bead-based size selection is essential. Always check the
library size distribution on a Bioanalyzer/TapeStation before sequencing.
\end{warningbox}

\subsection{miRNA-seq Analysis Workflow}

The analysis of miRNA-seq data follows a specialized pipeline:

\begin{enumerate}[nosep]
  \item \textbf{Adapter trimming:} Critical because inserts ($\sim$22\,nt)
    are shorter than reads (typically 50--75\,bp). Tools: \software{Cutadapt},
    \software{Trim Galore}, \software{fastp}.
  \item \textbf{Quality filtering:} Remove reads $<$ 15\,nt or $>$ 35\,nt after
    trimming.
  \item \textbf{Mapping to miRNA reference:} Align to mature and precursor miRNA
    sequences from miRBase using \software{Bowtie} (allowing 0--1 mismatches)
    or dedicated tools.
  \item \textbf{Quantification:} Count reads mapping to each miRNA.
  \item \textbf{Normalization and differential expression:} Use
    \software{DESeq2} or \software{edgeR} as for mRNA, but with miRNA-specific
    considerations (lower library complexity, more discrete count
    distributions).
  \item \textbf{Target prediction:} \software{TargetScan}, \software{miRDB},
    \software{DIANA-microT}.
\end{enumerate}

Key tools for small RNA analysis include:
\begin{itemize}[nosep]
  \item \textbf{miRDeep2:} Discovers known and novel miRNAs by modeling the
    miRNA biogenesis pathway (Drosha/Dicer processing signatures).
  \item \textbf{sRNAbench:} Comprehensive small RNA analysis including miRNA,
    piRNA, tRFs, and isomiR detection.
  \item \textbf{miRge 3.0:} Fast miRNA alignment and quantification using a
    Bayesian approach.
  \item \textbf{Oasis 2.0:} Web-based platform for small RNA-seq analysis.
\end{itemize}


% ---------------------------------------------------------------------------
\section{Long-Read RNA Sequencing}
\label{sec:rnaseq:longread}

Short-read RNA-seq excels at quantification but struggles with isoform
resolution. When multiple transcript isoforms share exons, short reads
($\sim$150\,bp) cannot span the full transcript, making it impossible to
determine which combination of exons belongs to the same molecule. Long-read
RNA sequencing solves this by reading full-length transcripts.

\subsection{PacBio Iso-Seq}

The \textbf{Isoform Sequencing (Iso-Seq)} protocol from Pacific Biosciences
produces full-length transcript sequences without the need for assembly:

\begin{enumerate}[nosep]
  \item Reverse-transcribe poly-A$^+$ RNA into full-length cDNA using a
    template-switching oligo.
  \item Size-select cDNA (optional; can fractionate to enrich specific size
    ranges).
  \item Generate SMRTbell libraries and sequence on the PacBio Revio or Sequel
    IIe platform.
  \item Use the \software{IsoSeq3} pipeline to produce high-quality (HQ)
    full-length transcript isoforms with $>$99\% accuracy using CCS reads.
\end{enumerate}

Typical output: 5--15 million HiFi reads per SMRT Cell on Revio, yielding
hundreds of thousands of unique transcript isoforms.

\subsection{Oxford Nanopore RNA Sequencing}

Oxford Nanopore offers two RNA-seq approaches:

\begin{description}[style=nextline]
  \item[Direct RNA sequencing (RNA004)] Sequences native RNA molecules directly
    through the nanopore without reverse transcription or PCR. This preserves
    RNA modifications (m$^6$A, pseudouridine, m$^5$C) and poly-A tail length
    as part of the raw signal. Requires $\sim$500\,ng poly-A$^+$ RNA.
    Currently lower throughput ($\sim$1--4M reads per flow cell).
  \item[cDNA sequencing (SQK-PCS114)] Reverse-transcribes RNA into cDNA with
    template-switching, then sequences the cDNA. Higher throughput than direct
    RNA; compatible with PCR amplification for lower inputs.
\end{description}

\begin{keyconcept}[title={Advantages of Long-Read RNA-seq}]
\begin{itemize}[nosep]
  \item \textbf{True isoform identification:} Each read represents a single
    full-length transcript, eliminating assembly artifacts.
  \item \textbf{No computational assembly:} No need for transcript assembly
    tools that can generate chimeric isoforms.
  \item \textbf{RNA modification detection:} ONT direct RNA-seq detects
    base modifications in native RNA.
  \item \textbf{Poly-A tail length measurement:} Direct information on
    polyadenylation dynamics.
  \item \textbf{Fusion transcript detection:} Full-length reads span fusion
    junctions unambiguously.
\end{itemize}
\end{keyconcept}

\subsection{Long-Read Isoform Analysis Tools}

\begin{description}[style=nextline]
  \item[\software{IsoQuant}] Maps long reads to a reference genome and assigns
    them to annotated or novel isoforms. Handles both PacBio and ONT data with
    sophisticated error correction.
  \item[\software{FLAIR}] Full-Length Alternative Isoform analysis of RNA.
    Corrects, collapses, and quantifies long-read isoforms. Integrates with
    short-read junction data for hybrid analysis.
  \item[\software{Bambu}] Reference-guided transcript discovery and
    quantification for long-read data. Uses a machine-learning model to
    distinguish genuine novel isoforms from noise.
  \item[\software{TALON}] Technology-agnostic long-read analysis pipeline;
    builds a database of observed transcripts across samples and tracks
    isoform usage changes.
  \item[\software{SQANTI3}] Structural and quality annotation of novel
    transcript isoforms; classifies isoforms into categories (full splice
    match, novel in catalog, novel not in catalog, etc.) and filters
    artifacts.
\end{description}


% ---------------------------------------------------------------------------
\section{Nascent RNA Sequencing}
\label{sec:rnaseq:nascent}

Standard RNA-seq measures \emph{steady-state} RNA levels, which reflect the
balance of transcription and degradation. Nascent RNA-seq methods specifically
capture RNA that is \emph{currently being transcribed} by RNA polymerase II,
providing a direct readout of transcriptional activity.

\begin{description}[style=nextline]
  \item[GRO-seq (Global Run-On sequencing)] Nuclei are isolated and RNA
    polymerases are allowed to extend nascent transcripts in the presence of
    BrUTP (bromouridine), which labels newly synthesized RNA. Labeled RNA is
    immunoprecipitated and sequenced. Reveals active transcription sites and
    their directionality, including enhancer RNAs (eRNAs).
  \item[PRO-seq (Precision Run-On sequencing)] Similar to GRO-seq but uses
    biotin-labeled NTPs that cause chain termination after a single nucleotide
    incorporation. This maps the position of active RNA polymerase at
    single-nucleotide resolution.
  \item[NET-seq (Native Elongating Transcript sequencing)] Immunoprecipitates
    RNA Pol II and sequences the 3$'$ ends of associated nascent RNA. No
    run-on step; captures the native polymerase position.
  \item[TT-seq (Transient Transcriptome sequencing)] Cells are pulsed with
    4-thiouridine (4sU) for 5--10 minutes; newly synthesized RNA incorporates
    4sU, which is biotinylated and selected. Captures both stable and unstable
    nascent transcripts at near-physiological conditions.
\end{description}

\begin{tipbox}
Nascent RNA-seq is essential for studying transcriptional regulation, enhancer
activity, and RNA polymerase dynamics. If your biological question is ``how much
RNA is being produced?'' rather than ``how much RNA is present?'', consider a
nascent RNA method.
\end{tipbox}


% ---------------------------------------------------------------------------
\section{RNA-seq Bioinformatics Workflow}
\label{sec:rnaseq:bioinformatics}

% TikZ: RNA-seq analysis workflow
\begin{figure}[htbp]
\centering
\begin{tikzpicture}[
  node distance=0.8cm and 1.5cm,
  every node/.style={font=\small},
  boxstep/.style={draw, rounded corners, fill=MidnightBlue!12, minimum width=3.8cm, minimum height=0.8cm, align=center, font=\small},
  tool/.style={font=\scriptsize\ttfamily, text=OliveGreen!80!black},
  arrow/.style={-{Stealth[length=5pt]}, thick, MidnightBlue!70},
  branch/.style={-{Stealth[length=5pt]}, thick, BrickRed!70},
]

% Main workflow
\node[boxstep] (raw) {Raw FASTQ reads};
\node[boxstep, below=of raw] (qc) {Quality control};
\node[boxstep, below=of qc] (trim) {Adapter/quality trimming};

% Branch point
\node[boxstep, below left=1.2cm and 1.0cm of trim] (align) {Splice-aware\\alignment};
\node[boxstep, below right=1.2cm and 1.0cm of trim] (pseudo) {Pseudo-alignment /\\lightweight quantification};

% After alignment
\node[boxstep, below=of align] (count) {Gene/transcript\\quantification};

% Merge
\node[boxstep, below=2.8cm of trim] (de) {Differential expression\\analysis};

% Downstream
\node[boxstep, below left=1.0cm and 0.8cm of de] (gsea) {Gene set\\enrichment};
\node[boxstep, below=of de] (iso) {Isoform / splicing\\analysis};
\node[boxstep, below right=1.0cm and 0.8cm of de] (fusion) {Fusion gene\\detection};

% Arrows
\draw[arrow] (raw) -- (qc);
\draw[arrow] (qc) -- (trim);
\draw[arrow] (trim) -| (align);
\draw[arrow] (trim) -| (pseudo);
\draw[arrow] (align) -- (count);
\draw[arrow] (count) -- (de);
\draw[arrow] (pseudo) |- (de);
\draw[arrow] (de) -| (gsea);
\draw[arrow] (de) -- (iso);
\draw[arrow] (de) -| (fusion);

% Tool labels
\node[tool, right=0.1cm of qc] {FastQC, MultiQC};
\node[tool, right=0.1cm of trim] {fastp, Trim Galore};
\node[tool, left=0.1cm of align] {STAR, HISAT2};
\node[tool, right=0.1cm of pseudo] {Salmon, kallisto};
\node[tool, left=0.1cm of count] {featureCounts, RSEM};
\node[tool, right=0.1cm of de] {DESeq2, edgeR};
\node[tool, left=0.05cm of gsea] {fgsea, clusterProfiler};
\node[tool, right=0.05cm of fusion] {STAR-Fusion, Arriba};
\node[tool, below=0.05cm of iso] {rMATS, SUPPA2, Leafcutter};

\end{tikzpicture}
\caption{Overview of the bulk RNA-seq bioinformatics workflow. Two main paths
exist after trimming: alignment-based quantification (left, using STAR/HISAT2
followed by featureCounts) or pseudo-alignment-based quantification (right,
using Salmon/kallisto). Both converge at differential expression analysis.}
\label{fig:rnaseq:workflow}
\end{figure}

\subsection{Quality Control}

\begin{description}[style=nextline]
  \item[\software{FastQC}] Generates per-read-base quality scores, GC content
    distribution, adapter contamination levels, sequence duplication levels,
    and overrepresented sequences for each FASTQ file.
  \item[\software{MultiQC}] Aggregates FastQC (and other tool) reports across
    all samples into a single interactive HTML report. Essential for spotting
    batch effects and outlier samples.
\end{description}

\textbf{Key RNA-seq QC metrics} to monitor:
\begin{itemize}[nosep]
  \item Per-base sequence quality (Phred $\geq$ 30 across the read).
  \item GC content distribution (should match species expectation; bimodal
    distribution suggests contamination).
  \item Adapter content (should be minimal for RNA-seq with appropriate insert
    sizes).
  \item Duplication rate (high duplication may indicate low-complexity library
    or over-amplification).
  \item After alignment: percentage of uniquely mapped reads, percentage of
    reads in exons vs. introns vs. intergenic regions, 5$'$-to-3$'$ gene body
    coverage (detects RNA degradation bias).
\end{itemize}

\subsection{Trimming}

\begin{description}[style=nextline]
  \item[\software{fastp}] All-in-one tool for adapter trimming, quality
    trimming, polyG tail removal (common on NovaSeq two-color chemistry),
    and QC report generation. Fast (multi-threaded C++).
  \item[\software{Trim Galore}] Wrapper around \software{Cutadapt} with
    automatic adapter detection and optional FastQC integration. Widely used in
    RNA-seq pipelines.
\end{description}

\begin{tipbox}
For modern Illumina data, \software{fastp} is often sufficient as the sole
preprocessing step. It performs adapter trimming, quality trimming, and
generates a QC report in a single pass, simplifying the workflow.
\end{tipbox}

\subsection{Alignment: Splice-Aware Aligners}

RNA-seq reads frequently span exon--exon junctions. Standard DNA aligners like
BWA cannot handle these split reads because they expect contiguous alignment
to the reference.

\begin{warningbox}
Never use \software{BWA} or \software{Bowtie2} (DNA aligners) for RNA-seq
read alignment. These tools are not splice-aware and will discard or misalign
reads that span exon--exon junctions, which can constitute 20--40\% of all
reads.
\end{warningbox}

\begin{description}[style=nextline]
  \item[\software{STAR} (Spliced Transcripts Alignment to a Reference)]
    The most widely used RNA-seq aligner. Uses a suffix array-based algorithm
    that is extremely fast but memory-intensive ($\sim$30--40\,GB RAM for
    human genome). Features:
    \begin{itemize}[nosep]
      \item Two-pass mode for improved novel junction detection.
      \item Built-in gene quantification (\texttt{--quantMode GeneCounts}).
      \item Chimeric read detection for fusion gene calling.
      \item \software{STARsolo} mode for single-cell RNA-seq (see \cref{ch:scrnaseq}).
    \end{itemize}
  \item[\software{HISAT2} (Hierarchical Indexing for Spliced Alignment of Transcripts)]
    Uses a graph-based FM index that is much more memory-efficient ($\sim$8\,GB
    for human). Slightly slower than STAR but suitable for machines with limited
    RAM. Successor to TopHat2/Bowtie2.
\end{description}

\subsection{Pseudo-Alignment and Lightweight Quantification}

An alternative to alignment-based approaches, pseudo-alignment (also called
quasi-mapping or lightweight alignment) estimates transcript abundance without
producing a traditional BAM file:

\begin{description}[style=nextline]
  \item[\software{Salmon}] Maps reads to the transcriptome (not genome) using
    selective alignment or quasi-mapping. Corrects for GC bias, sequence-specific
    bias, and fragment length distribution. Outputs transcript-level
    quantification (TPM and estimated counts). Run time: minutes rather than
    hours.
  \item[\software{kallisto}] Uses pseudo-alignment based on k-mer hashing.
    Produces transcript-level abundance estimates. Very fast and
    memory-efficient.
\end{description}

\begin{keyconcept}[title={Alignment vs. Pseudo-Alignment}]
\textbf{Alignment-based} (STAR $\rightarrow$ featureCounts): Produces BAM files
useful for visualization, variant calling, and splicing analysis. Required when
you need genome-level information. \textbf{Pseudo-alignment} (Salmon/kallisto):
Faster, produces transcript-level quantification directly. Excellent for
differential expression and isoform quantification. No BAM file for
visualization. For most differential expression analyses, both approaches yield
highly concordant results.
\end{keyconcept}

\subsection{Gene and Transcript Quantification}

\begin{description}[style=nextline]
  \item[\software{featureCounts} (Subread package)] Assigns aligned reads (from
    BAM files) to genomic features (genes, exons) based on a GTF/GFF
    annotation. Fast and simple. Produces a gene-by-sample count matrix.
    Handles paired-end reads, multi-mapping reads, and strand-specific
    assignments.
  \item[\software{RSEM}] Reference-free estimation of transcript-level
    expression using an expectation-maximization (EM) algorithm. Often used with
    STAR alignment (\texttt{STAR --quantMode TranscriptomeSAM}).
  \item[\software{StringTie}] Assembles transcripts from aligned reads and
    estimates abundance. Particularly useful for novel transcript discovery and
    isoform-level quantification.
\end{description}

\subsection{Differential Expression Analysis}

Differential expression (DE) analysis identifies genes whose expression differs
significantly between experimental conditions. The three dominant tools all use
negative binomial generalized linear models:

\begin{description}[style=nextline]
  \item[\software{DESeq2} (R/Bioconductor)] Uses shrinkage estimation for
    dispersion and fold changes. Provides independent filtering and Benjamini--Hochberg
    adjusted $p$-values. The most cited DE tool. Recommended for studies with
    $\geq$ 3 replicates per condition.
  \item[\software{edgeR} (R/Bioconductor)] Uses empirical Bayes methods for
    dispersion estimation. Slightly more liberal than DESeq2 in some benchmarks.
    Offers both exact tests and GLM-based approaches (quasi-likelihood F-tests
    are recommended for complex designs).
  \item[\software{limma-voom} (R/Bioconductor)] Transforms count data using the
    voom method (variance modeling at the observational level) and applies the
    limma linear modeling framework. Particularly powerful for complex
    experimental designs with multiple factors and interactions. Often performs
    well with larger sample sizes ($n > 10$ per group).
\end{description}

\begin{tipbox}
When importing Salmon or kallisto transcript-level counts into R, use the
\software{tximport} or \software{tximeta} package. These tools correctly
aggregate transcript-level estimates to gene-level counts while accounting for
changes in transcript length across samples---a bias that simple summing ignores.
\end{tipbox}

\subsection{The Negative Binomial Distribution for RNA-seq}
\label{subsec:rnaseq:negbinom}

RNA-seq read counts are \emph{discrete} non-negative integers, making continuous
distributions like the normal distribution inappropriate. A natural first model
is the Poisson distribution, where the variance equals the mean:
$\text{Var}(Y) = \mu$. However, biological replicates exhibit
\textbf{overdispersion}---the observed variance exceeds the Poisson expectation
because biological samples differ from one another beyond what sampling noise
alone would predict. A gene with mean expression $\mu = 100$ across replicates
may have variance of 500 or more, far larger than the Poisson prediction of 100.

The \textbf{negative binomial (NB) distribution} accounts for this
overdispersion by introducing a dispersion parameter $\alpha$ (also written
$\phi$ or $1/r$ in different parameterizations):
\begin{equation}
  \label{eq:rnaseq:nbvar}
  \text{Var}(Y) = \mu + \alpha\mu^2
\end{equation}
where $\mu$ is the mean expression and $\alpha \geq 0$ is the gene-specific
dispersion. When $\alpha = 0$, the NB reduces to the Poisson. For most genes in
a typical RNA-seq experiment, $\alpha$ ranges from 0.01 to 0.5. The probability
mass function of the NB distribution (using the mean-dispersion
parameterization) is:
\begin{equation}
  \label{eq:rnaseq:nbpmf}
  P(Y = k) = \binom{k + r - 1}{k}
    \left(\frac{r}{r + \mu}\right)^{r}
    \left(\frac{\mu}{r + \mu}\right)^{k},
    \qquad r = \frac{1}{\alpha}
\end{equation}
where $r = 1/\alpha$ is the size (or inverse dispersion) parameter.

\begin{keyconcept}[title={Why Not Poisson?}]
The Poisson model underestimates variance in RNA-seq data because it cannot
represent biological variability between replicates. Using a Poisson test
(e.g., an exact Poisson test in early RNA-seq tools) produces dramatically
inflated false-positive rates. All modern DE tools---\software{DESeq2},
\software{edgeR}, and \software{limma-voom}---use models that accommodate
overdispersion: the negative binomial GLM (\software{DESeq2}, \software{edgeR})
or variance-weighted linear models (\software{limma-voom}).
\end{keyconcept}

\subsection{DESeq2 Dispersion Estimation and Shrinkage}
\label{subsec:rnaseq:deseq2dispersion}

Accurate estimation of the gene-specific dispersion $\alpha_i$ is the central
statistical challenge in RNA-seq differential expression. With only 3--6
replicates per condition, per-gene maximum-likelihood estimates of $\alpha_i$
are highly noisy. \software{DESeq2} addresses this through a three-step
\textbf{dispersion shrinkage} procedure:

\begin{enumerate}[nosep]
  \item \textbf{Gene-wise dispersion estimates.} For each gene $i$,
    \software{DESeq2} computes a maximum-likelihood estimate $\hat{\alpha}_i^{\text{MLE}}$
    by fitting a negative binomial GLM using the Cox--Reid adjusted profile
    likelihood. These estimates are highly variable when the number of replicates
    is small.
  \item \textbf{Fitted dispersion trend.} A smooth function
    $\alpha_{\text{tr}}(\mu)$ is fitted to the gene-wise estimates as a function
    of the mean expression $\mu_i$. This trend captures the systematic
    relationship between dispersion and expression level (lowly expressed genes
    tend to have higher dispersion). The parametric fit is typically of the form:
    \begin{equation}
      \label{eq:rnaseq:disptrend}
      \alpha_{\text{tr}}(\mu) = \frac{a_1}{\mu} + a_0
    \end{equation}
    where $a_0$ and $a_1$ are estimated by regression.
  \item \textbf{Shrinkage toward the trend.} Each gene-wise estimate is
    shrunk toward the trend using a log-normal prior centered on
    $\log(\alpha_{\text{tr}}(\mu_i))$. The final maximum a posteriori (MAP)
    estimate balances the gene-specific data with the global trend:
    \begin{equation}
      \label{eq:rnaseq:dispmap}
      \hat{\alpha}_i^{\text{MAP}} = \arg\max_{\alpha_i}
        \left[
          \ell(\alpha_i \mid Y_i) +
          \log \pi(\alpha_i \mid \alpha_{\text{tr}}(\mu_i), \sigma_d^2)
        \right]
    \end{equation}
    where $\ell$ is the log-likelihood and $\pi$ is the log-normal prior with
    variance $\sigma_d^2$ estimated from the data. Genes with very high
    dispersion (outliers above the trend) are not shrunk, to avoid
    underestimating their variance and inflating false positives.
\end{enumerate}

\begin{tipbox}
The \software{edgeR} package uses a conceptually similar but technically
different approach. It estimates a common dispersion across all genes and
gene-wise dispersions, then shrinks gene-wise estimates toward the common
(or trended) dispersion using a weighted likelihood empirical Bayes procedure.
The key difference is that \software{edgeR} uses a weighted likelihood approach
rather than a prior distribution, and its quasi-likelihood F-test (recommended
for complex designs) adds a second layer of overdispersion modeling.
\end{tipbox}

\subsection{Normalisation Methods for RNA-seq}
\label{subsec:rnaseq:normalisation}

Raw read counts cannot be compared directly across genes or samples because they
are confounded by gene length, library size (total sequencing depth), and RNA
composition. Several normalisation metrics address different subsets of these
biases:

\begin{description}[style=nextline]
  \item[RPKM (Reads Per Kilobase per Million mapped reads)]
    Normalises for both gene length and sequencing depth. Used for single-end
    data:
    \begin{equation}
      \label{eq:rnaseq:rpkm}
      \text{RPKM}_i = \frac{c_i \times 10^9}{N \times l_i}
    \end{equation}
    where $c_i$ is the raw count for gene $i$, $N$ is the total number of mapped
    reads in the sample, and $l_i$ is the effective gene length in base pairs.
  \item[FPKM (Fragments Per Kilobase per Million mapped fragments)]
    The paired-end equivalent of RPKM, where $N$ counts fragments (read pairs)
    rather than individual reads. The formula is identical to RPKM but uses
    fragment counts, avoiding double-counting of paired-end reads.
  \item[TPM (Transcripts Per Million)]
    A two-step normalisation that first divides by gene length, then normalises
    to a fixed sum:
    \begin{equation}
      \label{eq:rnaseq:tpm}
      \text{TPM}_i = \frac{c_i / l_i}{\sum_{j} (c_j / l_j)} \times 10^6
    \end{equation}
    Because TPM values within each sample sum to exactly $10^6$, the relative
    proportions are directly comparable across samples---a property that RPKM/FPKM
    lack (their column sums vary with the gene length distribution of expressed
    genes).
\end{description}

\begin{warningbox}
RPKM and FPKM values are \textbf{not directly comparable across samples}
because the denominator $N$ includes all mapped reads, and the gene-length
normalisation step alters the effective library size differently depending on
which genes are expressed in each sample. TPM resolves this by normalising to a
fixed sum. For this reason, TPM has largely replaced RPKM/FPKM in modern
publications. However, for differential expression analysis, none of these
within-sample metrics are appropriate---use the raw counts with \software{DESeq2}
or \software{edgeR}, which apply their own between-sample normalisation
(median-of-ratios or TMM, respectively).
\end{warningbox}

\begin{keyconcept}[title={DESeq2 Median-of-Ratios Normalisation}]
\software{DESeq2} computes a \emph{size factor} $s_j$ for each sample $j$ as
the median of the ratios of observed counts to the geometric mean across
samples:
\begin{equation}
  \label{eq:rnaseq:sizefactor}
  s_j = \text{median}_i \frac{c_{ij}}{\left(\prod_{j'=1}^{m} c_{ij'}\right)^{1/m}}
\end{equation}
This approach is robust to highly expressed genes that dominate a few samples
and accounts for RNA composition differences---a major advantage over simple
library-size normalisation.
\end{keyconcept}

\subsection{DESeq2 vs.\ edgeR vs.\ limma-voom: A Detailed Comparison}
\label{subsec:rnaseq:detools_comparison}

\begin{table}[htbp]
\centering
\caption{Detailed comparison of the three major differential expression tools
for bulk RNA-seq.}
\label{tab:rnaseq:de_comparison}
\small
\begin{tabularx}{\textwidth}{lXXX}
\toprule
\textbf{Feature} & \textbf{DESeq2} & \textbf{edgeR} & \textbf{limma-voom} \\
\midrule
Statistical model & Negative binomial GLM & Negative binomial GLM (QL F-test recommended) & Weighted linear model on log-CPM \\
Normalisation & Median-of-ratios (size factors) & Trimmed mean of M-values (TMM) & TMM (via edgeR) or quantile \\
Dispersion estimation & Gene-wise MLE $\rightarrow$ trend $\rightarrow$ MAP shrinkage (log-normal prior) & Gene-wise $\rightarrow$ trended $\rightarrow$ empirical Bayes squeeze (weighted likelihood) & Mean-variance trend modeled by voom; precision weights per observation \\
Shrinkage of fold changes & Optional (apeglm, ashr, normal) & No built-in LFC shrinkage & Empirical Bayes moderation of $t$-statistics \\
Multiple testing & Benjamini--Hochberg with independent filtering & Benjamini--Hochberg & Benjamini--Hochberg \\
Complex designs & Supported (interaction terms, continuous covariates) & Excellent (quasi-likelihood framework is very flexible) & Excellent (full linear model framework; best for large, complex designs) \\
Speed (20k genes, $n$=20) & $\sim$1--3 minutes & $\sim$30 seconds--1 minute & $\sim$10--30 seconds \\
Best for & Standard two-group comparisons with small $n$; strong defaults & Complex experimental designs; large datasets & Very large sample sizes ($n > 20$); complex multi-factor designs \\
\bottomrule
\end{tabularx}
\end{table}

\begin{tipbox}
In practice, all three tools produce highly overlapping results for well-designed
experiments with adequate replication. \software{DESeq2} is the most popular due
to its sensible defaults and extensive documentation. \software{edgeR} with
quasi-likelihood F-tests is preferred for complex designs with multiple factors.
\software{limma-voom} excels when sample sizes are large (e.g., GTEx-scale
studies with hundreds of samples) because its linear model framework is both
fast and statistically powerful. For the most robust results, consider running
two tools and focusing on the intersection of significant genes.
\end{tipbox}

\subsection{Gene Set Enrichment Analysis}

After identifying differentially expressed genes, the next step is biological
interpretation:

\begin{description}[style=nextline]
  \item[\software{GSEA} (Gene Set Enrichment Analysis)] The original
    Broad Institute tool. Uses a ranked list of all genes (not just significant
    ones) to determine whether predefined gene sets (e.g., GO terms, KEGG
    pathways, MSigDB collections) are enriched at the top or bottom of the
    ranking. Avoids the arbitrary cutoff problem of over-representation analysis.
  \item[\software{fgsea} (R/Bioconductor)] Fast implementation of the GSEA
    algorithm in R. Orders of magnitude faster than the original Java tool.
  \item[\software{clusterProfiler} (R/Bioconductor)] Comprehensive enrichment
    analysis package supporting Gene Ontology (GO), KEGG, Reactome, Disease
    Ontology, and custom gene sets. Performs both over-representation analysis
    (ORA) and GSEA. Provides publication-quality visualization functions
    (dot plots, enrichment maps, ridge plots).
  \item[\software{enrichR}] Web-based tool with an R API that queries $>$200
    gene set libraries including ENCODE, ChEA, and other curated databases.
\end{description}

\subsection{Isoform and Splicing Analysis}

\begin{description}[style=nextline]
  \item[\software{rMATS} (replicate Multivariate Analysis of Transcript Splicing)]
    Detects differential alternative splicing events from RNA-seq data.
    Classifies events into five types: skipped exon (SE), retained intron
    (RI), mutually exclusive exons (MXE), alternative 5$'$ splice site (A5SS),
    and alternative 3$'$ splice site (A3SS). Uses a Bayesian framework to
    estimate inclusion levels ($\Psi$ / PSI values) and their differences.
  \item[\software{SUPPA2}] Fast differential splicing analysis using percent
    spliced-in (PSI) values calculated directly from transcript quantification
    (Salmon/kallisto output). Does not require BAM files.
  \item[\software{Leafcutter}] Detects differential intron usage without
    requiring transcript annotations. Identifies intron excision clusters and
    tests for differential usage between conditions. Annotation-free approach
    is particularly valuable for non-model organisms.
\end{description}

\subsection{Fusion Gene Detection}

Gene fusions---where two normally separate genes are joined---are hallmarks of
many cancers (e.g., \gene{BCR-ABL1} in CML, \gene{EML4-ALK} in lung cancer).

\begin{description}[style=nextline]
  \item[\software{STAR-Fusion}] Uses chimeric read alignments from STAR to
    detect fusion transcripts. Integrates with the FusionInspector
    visualization tool and with CTAT (Cancer Transcriptome Analysis Toolkit).
  \item[\software{Arriba}] Fast and accurate fusion detection from STAR chimeric
    alignments. Known for low false-positive rate and built-in visualization
    of fusion events with gene structure diagrams.
\end{description}


% ---------------------------------------------------------------------------
\section{Complete Snakemake Workflow for RNA-seq}
\label{sec:rnaseq:snakemake}

\begin{lstlisting}[language=Python, caption={Complete Snakemake workflow for
bulk RNA-seq analysis with STAR alignment and featureCounts quantification.},
label={lst:rnaseq:snakemake}, float=htbp]
# Snakefile for Bulk RNA-seq Analysis
# ====================================
# Pipeline: FastQC -> fastp -> STAR -> featureCounts -> DESeq2

import pandas as pd

configfile: "config/config.yaml"

# Load sample sheet
samples = pd.read_csv(config["samples"], sep="\t")
SAMPLES = samples["sample"].tolist()

# Reference files
GENOME_DIR = config["star_index"]
GTF = config["gtf"]
FASTA = config["genome_fasta"]

rule all:
    input:
        "results/multiqc/multiqc_report.html",
        "results/counts/gene_counts.tsv",
        "results/deseq2/differential_expression.csv",
        expand("results/fastqc_raw/{sample}_{read}_fastqc.html",
               sample=SAMPLES, read=["R1", "R2"]),
        expand("results/bigwig/{sample}.bw", sample=SAMPLES)

rule fastqc_raw:
    input:
        "data/raw/{sample}_{read}.fastq.gz"
    output:
        html="results/fastqc_raw/{sample}_{read}_fastqc.html",
        zip="results/fastqc_raw/{sample}_{read}_fastqc.zip"
    threads: 2
    conda: "envs/qc.yaml"
    shell:
        "fastqc -t {threads} -o results/fastqc_raw/ {input}"

rule fastp:
    input:
        r1="data/raw/{sample}_R1.fastq.gz",
        r2="data/raw/{sample}_R2.fastq.gz"
    output:
        r1="results/trimmed/{sample}_R1.fastq.gz",
        r2="results/trimmed/{sample}_R2.fastq.gz",
        html="results/fastp/{sample}.html",
        json="results/fastp/{sample}.json"
    threads: 4
    conda: "envs/qc.yaml"
    shell:
        """
        fastp \
          --in1 {input.r1} --in2 {input.r2} \
          --out1 {output.r1} --out2 {output.r2} \
          --html {output.html} --json {output.json} \
          --detect_adapter_for_pe \
          --qualified_quality_phred 20 \
          --length_required 35 \
          --thread {threads}
        """

rule star_index:
    input:
        fasta=FASTA,
        gtf=GTF
    output:
        directory("resources/star_index")
    threads: 12
    resources:
        mem_mb=40000
    conda: "envs/align.yaml"
    shell:
        """
        STAR --runMode genomeGenerate \
          --genomeDir {output} \
          --genomeFastaFiles {input.fasta} \
          --sjdbGTFfile {input.gtf} \
          --sjdbOverhang 149 \
          --runThreadN {threads}
        """

rule star_align:
    input:
        r1="results/trimmed/{sample}_R1.fastq.gz",
        r2="results/trimmed/{sample}_R2.fastq.gz",
        index=GENOME_DIR
    output:
        bam="results/star/{sample}/Aligned.sortedByCoord.out.bam",
        log="results/star/{sample}/Log.final.out",
        sj="results/star/{sample}/SJ.out.tab"
    threads: 8
    resources:
        mem_mb=36000
    params:
        prefix="results/star/{sample}/"
    conda: "envs/align.yaml"
    shell:
        """
        STAR --runMode alignReads \
          --genomeDir {input.index} \
          --readFilesIn {input.r1} {input.r2} \
          --readFilesCommand zcat \
          --outSAMtype BAM SortedByCoordinate \
          --outSAMstrandField intronMotif \
          --quantMode GeneCounts \
          --twopassMode Basic \
          --outFileNamePrefix {params.prefix} \
          --runThreadN {threads} \
          --outSAMattributes NH HI AS NM MD \
          --outFilterMultimapNmax 20 \
          --alignSJoverhangMin 8 \
          --alignSJDBoverhangMin 1 \
          --chimSegmentMin 20 \
          --chimJunctionOverhangMin 20
        """

rule index_bam:
    input:
        "results/star/{sample}/Aligned.sortedByCoord.out.bam"
    output:
        "results/star/{sample}/Aligned.sortedByCoord.out.bam.bai"
    conda: "envs/align.yaml"
    shell:
        "samtools index {input}"

rule featurecounts:
    input:
        bams=expand(
            "results/star/{sample}/Aligned.sortedByCoord.out.bam",
            sample=SAMPLES),
        gtf=GTF
    output:
        counts="results/counts/gene_counts.tsv",
        summary="results/counts/gene_counts.tsv.summary"
    threads: 4
    params:
        strandedness=2  # 2 = reverse-stranded (dUTP)
    conda: "envs/quantify.yaml"
    shell:
        """
        featureCounts \
          -a {input.gtf} \
          -o {output.counts} \
          -p --countReadPairs \
          -s {params.strandedness} \
          -T {threads} \
          -B -C \
          --extraAttributes gene_name \
          {input.bams}
        """

rule bigwig:
    input:
        bam="results/star/{sample}/Aligned.sortedByCoord.out.bam",
        bai="results/star/{sample}/Aligned.sortedByCoord.out.bam.bai"
    output:
        "results/bigwig/{sample}.bw"
    threads: 4
    conda: "envs/quantify.yaml"
    shell:
        """
        bamCoverage -b {input.bam} -o {output} \
          --normalizeUsing RPKM \
          --binSize 10 \
          -p {threads}
        """

rule multiqc:
    input:
        expand("results/fastqc_raw/{sample}_{read}_fastqc.zip",
               sample=SAMPLES, read=["R1", "R2"]),
        expand("results/fastp/{sample}.json", sample=SAMPLES),
        expand("results/star/{sample}/Log.final.out",
               sample=SAMPLES),
        "results/counts/gene_counts.tsv.summary"
    output:
        "results/multiqc/multiqc_report.html"
    conda: "envs/qc.yaml"
    shell:
        "multiqc results/ -o results/multiqc/ --force"

rule deseq2:
    input:
        counts="results/counts/gene_counts.tsv",
        samples=config["samples"]
    output:
        results="results/deseq2/differential_expression.csv",
        ma_plot="results/deseq2/ma_plot.pdf",
        volcano="results/deseq2/volcano_plot.pdf",
        pca="results/deseq2/pca_plot.pdf"
    conda: "envs/deseq2.yaml"
    script:
        "scripts/run_deseq2.R"
\end{lstlisting}


% ---------------------------------------------------------------------------
\section{Complete Nextflow Workflow for RNA-seq}
\label{sec:rnaseq:nextflow}

\begin{lstlisting}[language=Java, caption={Complete Nextflow DSL2 workflow for
bulk RNA-seq analysis with STAR alignment and featureCounts quantification.},
label={lst:rnaseq:nextflow}, float=htbp]
#!/usr/bin/env nextflow
// RNA-seq Analysis Pipeline (Nextflow DSL2)
// ==========================================
// Pipeline: FastQC -> fastp -> STAR -> featureCounts -> DESeq2

nextflow.enable.dsl = 2

params.reads       = "data/raw/*_{R1,R2}.fastq.gz"
params.genome      = "resources/genome.fa"
params.gtf         = "resources/genes.gtf"
params.star_index  = null
params.outdir      = "results"
params.strandedness = 2  // 2 = reverse-stranded (dUTP)

// Print pipeline info
log.info """
=============================
 RNA-seq Analysis Pipeline
=============================
 reads        : ${params.reads}
 genome       : ${params.genome}
 GTF          : ${params.gtf}
 strandedness : ${params.strandedness}
 outdir       : ${params.outdir}
=============================
"""

// --- Processes ---

process FASTQC {
    tag "$sample_id"
    publishDir "${params.outdir}/fastqc", mode: 'copy'
    cpus 2

    input:
    tuple val(sample_id), path(reads)

    output:
    path("*.html"), emit: html
    path("*.zip"),  emit: zip

    script:
    """
    fastqc -t ${task.cpus} ${reads}
    """
}

process FASTP {
    tag "$sample_id"
    publishDir "${params.outdir}/fastp", mode: 'copy'
    cpus 4

    input:
    tuple val(sample_id), path(reads)

    output:
    tuple val(sample_id), path("*_trimmed_R{1,2}.fastq.gz"),
        emit: trimmed_reads
    path("*.html"), emit: html
    path("*.json"), emit: json

    script:
    def (r1, r2) = reads
    """
    fastp \
      --in1 ${r1} --in2 ${r2} \
      --out1 ${sample_id}_trimmed_R1.fastq.gz \
      --out2 ${sample_id}_trimmed_R2.fastq.gz \
      --html ${sample_id}_fastp.html \
      --json ${sample_id}_fastp.json \
      --detect_adapter_for_pe \
      --qualified_quality_phred 20 \
      --length_required 35 \
      --thread ${task.cpus}
    """
}

process STAR_INDEX {
    publishDir "${params.outdir}/star_index", mode: 'copy'
    cpus 12
    memory '40 GB'

    input:
    path(genome)
    path(gtf)

    output:
    path("star_index"), emit: index

    script:
    """
    mkdir -p star_index
    STAR --runMode genomeGenerate \
      --genomeDir star_index \
      --genomeFastaFiles ${genome} \
      --sjdbGTFfile ${gtf} \
      --sjdbOverhang 149 \
      --runThreadN ${task.cpus}
    """
}

process STAR_ALIGN {
    tag "$sample_id"
    publishDir "${params.outdir}/star/${sample_id}", mode: 'copy'
    cpus 8
    memory '36 GB'

    input:
    tuple val(sample_id), path(reads)
    path(index)

    output:
    tuple val(sample_id),
        path("*Aligned.sortedByCoord.out.bam"), emit: bam
    path("*Log.final.out"),                     emit: log
    path("*SJ.out.tab"),                        emit: sj

    script:
    def (r1, r2) = reads
    """
    STAR --runMode alignReads \
      --genomeDir ${index} \
      --readFilesIn ${r1} ${r2} \
      --readFilesCommand zcat \
      --outSAMtype BAM SortedByCoordinate \
      --outSAMstrandField intronMotif \
      --quantMode GeneCounts \
      --twopassMode Basic \
      --outFileNamePrefix ${sample_id}_ \
      --runThreadN ${task.cpus} \
      --outSAMattributes NH HI AS NM MD \
      --outFilterMultimapNmax 20 \
      --chimSegmentMin 20
    """
}

process FEATURECOUNTS {
    publishDir "${params.outdir}/counts", mode: 'copy'
    cpus 4

    input:
    path(bams)
    path(gtf)

    output:
    path("gene_counts.tsv"),         emit: counts
    path("gene_counts.tsv.summary"), emit: summary

    script:
    """
    featureCounts \
      -a ${gtf} \
      -o gene_counts.tsv \
      -p --countReadPairs \
      -s ${params.strandedness} \
      -T ${task.cpus} \
      -B -C \
      ${bams}
    """
}

process MULTIQC {
    publishDir "${params.outdir}/multiqc", mode: 'copy'

    input:
    path('*')

    output:
    path("multiqc_report.html")

    script:
    """
    multiqc . -o . --force
    """
}

process DESEQ2 {
    publishDir "${params.outdir}/deseq2", mode: 'copy'

    input:
    path(counts)
    path(sample_sheet)

    output:
    path("differential_expression.csv"), emit: results
    path("*.pdf"), emit: plots

    script:
    """
    Rscript ${projectDir}/bin/run_deseq2.R \
      --counts ${counts} \
      --samples ${sample_sheet} \
      --outdir .
    """
}

// --- Workflow ---

workflow {
    // Create channel from read pairs
    reads_ch = Channel
        .fromFilePairs(params.reads, checkIfExists: true)

    // Quality control on raw reads
    FASTQC(reads_ch)

    // Adapter and quality trimming
    FASTP(reads_ch)

    // Build STAR index if not provided
    if (params.star_index) {
        star_index = Channel.fromPath(params.star_index)
    } else {
        STAR_INDEX(
            Channel.fromPath(params.genome),
            Channel.fromPath(params.gtf)
        )
        star_index = STAR_INDEX.out.index
    }

    // Align reads with STAR
    STAR_ALIGN(FASTP.out.trimmed_reads, star_index.collect())

    // Quantify with featureCounts
    FEATURECOUNTS(
        STAR_ALIGN.out.bam.map { it[1] }.collect(),
        Channel.fromPath(params.gtf)
    )

    // Aggregate QC reports
    MULTIQC(
        FASTQC.out.zip
            .mix(FASTP.out.json)
            .mix(STAR_ALIGN.out.log)
            .mix(FEATURECOUNTS.out.summary)
            .collect()
    )
}
\end{lstlisting}

\begin{tipbox}
For production use, consider the \software{nf-core/rnaseq} pipeline, which
implements a full-featured, peer-reviewed RNA-seq workflow in Nextflow with
extensive QC, multiple aligner/quantifier options, and compatibility with
institutional HPC systems and cloud environments. As of v3.14, it supports
STAR, HISAT2, Salmon, and RSEM, with automatic strandedness detection.
\end{tipbox}


% ---------------------------------------------------------------------------
\section{Experimental Design for RNA-seq}
\label{sec:rnaseq:design}

\subsection{Biological Replicates}

\begin{keyconcept}[title={Replication is Non-Negotiable}]
Biological replicates are essential for statistical inference. Technical
replicates (sequencing the same library twice) do not capture biological
variability and are rarely needed given the high reproducibility of modern
sequencing. The minimum is \textbf{3 biological replicates per condition};
\textbf{5--6 replicates} substantially increase power to detect smaller fold
changes. For heterogeneous clinical samples, even more replicates are
advisable.
\end{keyconcept}

Power analysis tools such as \software{RNASeqPower} (R/Bioconductor) and
\software{PROPER} can estimate the number of replicates needed given expected
effect sizes, dispersion, and sequencing depth.

\subsection{Sequencing Depth}

\begin{table}[htbp]
\centering
\caption{Recommended sequencing depths for common RNA-seq applications.}
\label{tab:rnaseq:depth}
\begin{tabularx}{\textwidth}{lXr}
\toprule
\textbf{Application} & \textbf{Rationale} & \textbf{Reads/sample} \\
\midrule
Differential gene expression & Sufficient to quantify $\sim$15,000 expressed genes & 20--30M \\
Isoform-level analysis & Need to distinguish splice variants & 50--100M \\
Novel transcript discovery & Deeper coverage of low-abundance transcripts & 80--150M \\
Small RNA-seq & Smaller library; many reads are adapter dimers & 10--20M \\
Total RNA-seq (rRNA depleted) & Compensate for residual rRNA reads & 40--60M \\
FFPE RNA-seq & Higher duplicate rate reduces usable reads & 50--80M \\
\bottomrule
\end{tabularx}
\end{table}

\begin{tipbox}
In most cases, adding biological replicates provides more statistical power
than increasing sequencing depth beyond 20--30M reads per sample. Budget your
sequencing funds toward more replicates rather than deeper sequencing, unless
you specifically need isoform-level resolution.
\end{tipbox}

\subsection{Batch Effects and Their Mitigation}

Batch effects are systematic technical differences between groups of samples
processed at different times, by different operators, or on different flow
cells. They can be devastating in RNA-seq because they may confound or entirely
obscure biological signal.

\textbf{Prevention (most important):}
\begin{itemize}[nosep]
  \item Extract RNA from all samples in the same session if possible.
  \item Prepare libraries in a single batch.
  \item If batching is unavoidable, \textbf{balance experimental groups across
    batches} (do not process all controls in batch 1 and all treatments in
    batch 2).
  \item Record batch information meticulously.
\end{itemize}

\textbf{Computational correction (second line of defense):}
\begin{description}[style=nextline]
  \item[\software{ComBat} (sva package)] Empirical Bayes framework for batch
    effect adjustment. Operates on normalized expression matrices. Widely used
    and effective but can remove real biological signal if batch and condition
    are confounded.
  \item[\software{limma::removeBatchEffect}] Removes batch effects from a
    normalized expression matrix for visualization purposes (PCA, heatmaps).
    For differential expression, include batch as a covariate in the DE model
    instead.
  \item[Inclusion in the DE model] The statistically preferred approach: add
    batch as a covariate in the DESeq2 or edgeR design formula (e.g.,
    \texttt{\textasciitilde{} batch + condition}). This separates batch
    variance from treatment variance without modifying the data.
\end{description}

\begin{warningbox}
If batch is completely confounded with the biological variable of interest
(e.g., all treatment samples were processed in January and all controls in
March), no computational method can rescue the experiment. The only solution
is to re-design and re-run the experiment with proper balancing.
\end{warningbox}


% ---------------------------------------------------------------------------
\section{RNA-seq Approaches: Comprehensive Comparison}
\label{sec:rnaseq:comparison}

\begin{table}[htbp]
\centering
\caption{Comprehensive comparison of RNA-seq approaches.}
\label{tab:rnaseq:approaches}
\small
\begin{tabularx}{\textwidth}{lXXXX}
\toprule
\textbf{Feature} & \textbf{mRNA-seq} & \textbf{Total RNA-seq} & \textbf{Small RNA-seq} & \textbf{Long-read RNA-seq} \\
\midrule
Target & Poly-A$^+$ mRNA & All RNA (rRNA depleted) & miRNA, piRNA, tRFs & Full-length transcripts \\
Read length & 2$\times$150\,bp & 2$\times$150\,bp & 1$\times$50--75\,bp & 1--20\,kb \\
Platform & Illumina & Illumina & Illumina & PacBio, ONT \\
Input RNA & 100\,ng--1\,$\mu$g & 100\,ng--1\,$\mu$g & 100\,ng--1\,$\mu$g & 300\,ng--5\,$\mu$g \\
Quality req. & RIN $\geq$ 7 & RIN $\geq$ 2 (with rRNA depl.) & RIN $\geq$ 7 & RIN $\geq$ 7 (Iso-Seq) \\
Isoform resolution & Limited & Limited & N/A & Excellent \\
ncRNA coverage & Minimal & Comprehensive & Small ncRNA only & Good \\
Depth needed & 20--30M & 40--60M & 10--20M & 1--10M \\
Cost/sample & \$\$  & \$\$ & \$ & \$\$\$ \\
Main use case & Standard DE & Degraded samples, ncRNA & miRNA profiling & Isoform discovery \\
\bottomrule
\end{tabularx}
\end{table}


% ---------------------------------------------------------------------------
\section{Chapter Summary}
\label{sec:rnaseq:summary}

RNA-seq is the workhorse technology of modern transcriptomics, applicable to
an extraordinary range of biological questions. In this chapter we covered:

\begin{itemize}[nosep]
  \item The rationale for measuring RNA and the complexity of the transcriptome.
  \item Poly-A selection (mRNA-seq) versus ribosomal depletion (total RNA-seq)
    and when to use each.
  \item Library preparation details: strand-specific protocols (dUTP method),
    low-input protocols (Smart-seq), and degraded RNA handling (FFPE).
  \item Small RNA sequencing for miRNAs and other small non-coding RNAs.
  \item Long-read RNA-seq (PacBio Iso-Seq, ONT direct RNA) for full-length
    isoform characterization.
  \item Nascent RNA methods (GRO-seq, PRO-seq, NET-seq, TT-seq) for measuring
    active transcription.
  \item The complete bioinformatics pipeline: QC, trimming, alignment (STAR,
    HISAT2), quantification (featureCounts, Salmon), differential expression
    (DESeq2, edgeR, limma-voom), gene set enrichment, isoform analysis, and
    fusion detection.
  \item Practical workflows in both Snakemake and Nextflow.
  \item Experimental design: replication, sequencing depth, and batch effect
    management.
\end{itemize}

In the next chapter, we move from bulk tissue-level measurements to the
resolution of individual cells with single-cell RNA sequencing
(\cref{ch:scrnaseq}).
